\documentclass[11pt,a4paper]{article}

\usepackage{amsmath, amssymb, amsthm}
\usepackage{geometry}
\usepackage{hyperref}
\usepackage{booktabs}
\usepackage{array}
\usepackage{xcolor}
\usepackage{parskip}
\usepackage{titlesec}
\usepackage{fancyhdr}
\usepackage{microtype}
\usepackage{bm}

\geometry{margin=2.5cm}

\hypersetup{
  colorlinks=true,
  linkcolor=blue!60!black,
  urlcolor=blue!60!black,
}

\pagestyle{fancy}
\fancyhf{}
\rhead{\small CADDack GNN Mathematics}
\lhead{\small \nouppercase{\leftmark}}
\cfoot{\thepage}

\titleformat{\section}{\large\bfseries}{{\S}\,\thesection}{1em}{}[\vspace{-4pt}\rule{\linewidth}{0.4pt}]
\titleformat{\subsection}{\normalsize\bfseries}{\thesubsection}{1em}{}

\newcommand{\R}{\mathbb{R}}
\newcommand{\N}{\mathcal{N}}
\newcommand{\bx}{\mathbf{x}}
\newcommand{\bh}{\mathbf{h}}
\newcommand{\be}{\mathbf{e}}
\newcommand{\br}{\mathbf{r}}
\newcommand{\bmu}{\boldsymbol{\mu}}
\newcommand{\bs}{\boldsymbol{\sigma}}
\newcommand{\bphi}{\boldsymbol{\phi}}
\newcommand{\softplus}{\mathrm{softplus}}
\newcommand{\ReLU}{\mathrm{ReLU}}
\newcommand{\SiLU}{\mathrm{SiLU}}
\newcommand{\MLP}{\mathrm{MLP}}
\newcommand{\code}[1]{\texttt{#1}}

\title{%
  \textbf{CADDack: GNN Mathematics Reference}\\[6pt]
  \large Two-Tower Bayesian Fusion Model for Protein--Ligand Binding Affinity
}
\author{CADDack Project}
\date{\today}

\begin{document}

\maketitle
\thispagestyle{fancy}

\begin{abstract}
This document gives a self-contained mathematical description of every computation
performed by the CADDack graph neural network stack, from raw molecular input to
predicted binding affinity with calibrated uncertainty.  The architecture consists of
(1) a 2D ligand message-passing encoder (GCN / GINE), (2) a 3D geometry tower
(SchNet-style continuous-filter convolution), (3) a Bayesian fusion head
(Bayes-by-Backprop, heteroscedastic ELBO), and (4) Monte Carlo uncertainty
decomposition at inference.  Each formula is cross-referenced to its implementing
source file.
\end{abstract}

\tableofcontents
\newpage

% ============================================================
\section{Input Representation}
\label{sec:input}
% ============================================================

\subsection{Atom feature vector (\code{datasets.py:\_atom\_features})}

Each heavy atom $i$ in a molecule is represented by a seven-dimensional feature
vector $\bx_i \in \R^7$:

\begin{equation}
\bx_i =
\bigl[
  z_i,\;
  \deg(i),\;
  q_i,\;
  \mathbf{1}_{\mathrm{arom}},\;
  \mathrm{hyb}_i,\;
  n^H_i,\;
  \mathbf{1}_{\mathrm{ring}}
\bigr]^{\!\top}
\end{equation}

where:
\begin{itemize}
  \item $z_i \in \mathbb{Z}_{>0}$ is the atomic number (e.g.\ $z=6$ for carbon, $z=7$ for nitrogen).
  \item $\deg(i)$ is the number of explicit bonds (degree in the molecular graph).
  \item $q_i$ is the formal charge.
  \item $\mathbf{1}_{\mathrm{arom}} \in \{0,1\}$ indicates aromaticity.
  \item $\mathrm{hyb}_i \in \{1,\ldots,6\}$ is the integer encoding of the hybridisation
        state (SP, SP2, SP3, \ldots).
  \item $n^H_i$ is the total hydrogen count (implicit + explicit).
  \item $\mathbf{1}_{\mathrm{ring}} \in \{0,1\}$ indicates ring membership.
\end{itemize}

All entries are cast to \texttt{float32}.

\subsection{Bond feature vector (\code{datasets.py:\_bond\_features})}

Each bond $(i,j)$ in the molecular graph is represented by a seven-dimensional edge
feature vector $\be_{ij} \in \R^7$:

\begin{equation}
\be_{ij} =
\bigl[
  \mathbf{1}_{\text{single}},\;
  \mathbf{1}_{\text{double}},\;
  \mathbf{1}_{\text{triple}},\;
  \mathbf{1}_{\text{arom}},\;
  \mathbf{1}_{\text{conj}},\;
  \mathbf{1}_{\text{ring}},\;
  \mathrm{stereo}
\bigr]^{\!\top}
\end{equation}

The first four elements form a one-hot encoding of bond order/type.
$\mathrm{stereo} \in \{0,1,2,\ldots\}$ is the integer cast of the RDKit stereo enum.

Edges are \textbf{bidirectional}: for each bond $(i,j)$ both directed edges
$(i \to j)$ and $(j \to i)$ are inserted with identical features.

% ============================================================
\section{Tower 1 --- Ligand 2D Message Passing}
\label{sec:tower1}
% ============================================================

\subsection{MolecularGCN (\code{models.py:MolecularGCN})}

The Graph Convolutional Network~\cite{kipf2017} applies the symmetric-normalised
spectral convolution at each layer $\ell$:

\begin{equation}
\bx_i^{(\ell+1)}
= \ReLU\!\left(
    W^{(\ell)} \sum_{j \in \mathcal{N}(i) \cup \{i\}}
    \frac{\bx_j^{(\ell)}}{\sqrt{d_i\, d_j}}
  \right)
\label{eq:gcn}
\end{equation}

where $d_i = \deg(i)+1$ (self-loop included) and $W^{(\ell)} \in \R^{C \times C}$ is
the learnable weight matrix.  Dropout with rate $p$ is applied after each
non-linearity.

After $L$ convolutional layers, global \emph{mean} pooling produces the molecular
embedding:

\begin{equation}
\bh_{\mathrm{mol}} = \frac{1}{|\mathcal{V}|} \sum_{i \in \mathcal{V}} \bx_i^{(L)}
\;\in \R^C
\end{equation}

The prediction head is a single affine map: $\hat{y} = \mathbf{w}^\top \bh_{\mathrm{mol}} + b$.

\subsection{MolecularGINE (\code{models.py:MolecularGINE})}

The Graph Isomorphism Network with Edge features (GINE)~\cite{xu2019,hu2020}
extends GIN~\cite{xu2019} to incorporate bond features.  At layer $\ell$:

\begin{equation}
\bx_i^{(\ell+1)}
= \MLP^{(\ell)}\!\left(
    (1 + \varepsilon)\,\bx_i^{(\ell)}
    + \sum_{j \in \mathcal{N}(i)} \ReLU\!\bigl(\bx_j^{(\ell)} + \be_{ij}\bigr)
  \right)
\label{eq:gine}
\end{equation}

Key design choices:
\begin{itemize}
  \item The edge features $\be_{ij}$ are added (not concatenated) to the sender
        embedding before aggregation.  This is algebraically equivalent to a
        first-order Taylor expansion of an edge-conditioned filter.
  \item Aggregation is \textbf{summation}, which is strictly more expressive than
        mean or max and makes GINE at least as powerful as the
        Weisfeiler--Leman (1-WL) graph isomorphism test~\cite{xu2019}.
  \item Each $\MLP^{(\ell)}$ is a two-layer network:
        $\operatorname{Linear}(C_{\mathrm{in}} \to C) \to \ReLU \to
        \operatorname{Linear}(C \to C)$.
\end{itemize}

After $L$ layers, global \emph{add} pooling:

\begin{equation}
\bh_{\mathrm{lig}} = \sum_{i \in \mathcal{V}} \bx_i^{(L)} \;\in \R^C
\end{equation}

GINE is the recommended ligand encoder because bond-order information (double bonds,
aromaticity) is essential for molecular property prediction.

% ============================================================
\section{Tower 2 --- 3D Geometry Tower}
\label{sec:tower2}
% ============================================================

Based on SchNet~\cite{schutt2017}.  The tower operates exclusively on
\emph{interatomic distances}, which guarantees $\mathrm{SE}(3)$ invariance.

\subsection{Atom embedding}

\begin{equation}
\bh_i^{(0)} = \mathrm{Embedding}(z_i) \;\in \R^C
\end{equation}

A learned look-up table maps each integer atomic number $z_i$ to a $C$-dimensional
vector.

\subsection{Radius graph}

Edges are created between all atom pairs within cutoff $c$ (default $c = 6.0$\,\AA):

\begin{equation}
\mathcal{E} = \bigl\{(i,j) : \|\br_i - \br_j\|_2 \leq c,\; i \neq j\bigr\}
\end{equation}

The edge set spans ligand--ligand, pocket--pocket, \emph{and} ligand--pocket pairs.
The ligand--pocket cross-edges represent the direct binding contacts that the model
must learn from.  The radius graph is computed \emph{inside} \code{forward()}
using the per-batch position tensors so that no pre-computed \code{edge\_index}
ever needs PyG collation across differently-sized graphs.

\subsection{Gaussian radial basis functions (\code{models.py:GaussianRBF})}

The scalar distance $d_{ij} = \|\br_i - \br_j\|_2$ is expanded into $K$ basis
functions:

\begin{equation}
\phi_k(d) = \exp\!\left( -\frac{(d - \mu_k)^2}{\gamma} \right), \quad k = 1, \ldots, K
\label{eq:rbf}
\end{equation}

where the centres $\mu_k$ are equally spaced on $[0, c]$:

\begin{equation}
\mu_k = \frac{(k-1)\,c}{K-1}, \qquad
\gamma = \left(\frac{c}{K}\right)^{\!2}
\end{equation}

The resulting feature vector $\bphi(d_{ij}) \in \R^K$ is a soft histogram over
distance ranges ($K = 50$ by default).

\subsection{Cosine envelope}

To ensure that messages vanish smoothly at the cutoff boundary (preventing
a discontinuity in the learned function at $d = c$):

\begin{equation}
w(d) = \tfrac{1}{2}\Bigl(\cos\!\Bigl(\frac{\pi\, d}{c}\Bigr) + 1\Bigr)
\label{eq:envelope}
\end{equation}

$w(0) = 1$, $w(c/2) = 1/2$, $w(c) = 0$.

\subsection{Interaction block (\code{models.py:\_InteractionBlock})}

Each of the $T$ interaction blocks ($T = 3$ by default) performs the following
steps.

\paragraph{Filter generation.}
\begin{equation}
\mathbf{f}_{ij}^{(\ell)}
= \mathrm{FilterNet}^{(\ell)}\!\bigl(\bphi(d_{ij})\bigr)
  \cdot w(d_{ij})
\;\in \R^C
\label{eq:filter}
\end{equation}

where $\mathrm{FilterNet}$ is a two-layer MLP with SiLU activations:

\begin{equation}
\mathrm{FilterNet}(\mathbf{u})
= W_2\,\SiLU(W_1\,\mathbf{u} + b_1) + b_2,
\qquad
\SiLU(x) = x\,\sigma(x)
\end{equation}

\paragraph{Continuous-filter message.}
\begin{equation}
\mathbf{m}_{ij}^{(\ell)}
= \bh_j^{(\ell)} \odot \mathbf{f}_{ij}^{(\ell)}
\;\in \R^C
\label{eq:cfconv}
\end{equation}

The sender embedding is modulated element-wise by the distance-dependent filter.
Unlike graph convolutions with fixed weight matrices, the filter here is a
\emph{continuous function of distance} that is learned end-to-end.

\paragraph{Aggregation and residual update.}
\begin{align}
\mathbf{a}_i^{(\ell)}
  &= \sum_{\substack{j:\,(i,j)\in\mathcal{E}}} \mathbf{m}_{ij}^{(\ell)}
\label{eq:agg}\\[4pt]
\bh_i^{(\ell+1)}
  &= \bh_i^{(\ell)} + \mathrm{UpdateMLP}^{(\ell)}\!\bigl(\mathbf{a}_i^{(\ell)}\bigr)
\label{eq:residual}
\end{align}

$\mathrm{UpdateMLP}$: $\operatorname{Linear}(C\to C) \to \SiLU \to
\operatorname{Linear}(C\to C)$.  The residual connection stabilises training.

\subsection{Readout}

\begin{equation}
\bh_{\mathrm{geo}} = \sum_{i \in \mathcal{V}} \bh_i^{(T)} \;\in \R^C
\end{equation}

\subsection{SE(3) invariance}

The geometry tower is invariant to all rigid-body transformations of the complex.
For any rotation $R \in \mathrm{SO}(3)$ and translation $\mathbf{t} \in \R^3$:

\begin{equation}
\bigl\|(R\,\br_i + \mathbf{t}) - (R\,\br_j + \mathbf{t})\bigr\|_2
= \|\br_i - \br_j\|_2
\end{equation}

Since every computation depends only on $d_{ij}$, the output is unchanged under
rotation or translation of the entire complex.

% ============================================================
\section{Bayesian Fusion Head}
\label{sec:bayes}
% ============================================================

Based on Bayes-by-Backprop~\cite{blundell2015}.  Towers are deterministic;
only the fusion head is Bayesian.  This keeps the KL tractable and training stable.

\subsection{Fusion}

The two tower outputs are concatenated:

\begin{equation}
\mathbf{h} = \bigl[\bh_{\mathrm{lig}};\, \bh_{\mathrm{geo}}\bigr] \in \R^{2C}
\end{equation}

\subsection{Variational posterior (\code{bayes.py:BayesianLinear})}

Each weight $w_{ij}$ in the fusion MLP has an independent Gaussian posterior:

\begin{equation}
q(w_{ij}) = \mathcal{N}\!\bigl(\mu_{ij},\, \sigma_{ij}^2\bigr)
\end{equation}

The standard deviation is parameterised via the softplus function to ensure
positivity:

\begin{equation}
\sigma_{ij} = \softplus(\rho_{ij}) = \ln\!\bigl(1 + e^{\rho_{ij}}\bigr)
\label{eq:softplus}
\end{equation}

Initialisation: $\mu_{ij}$ with Kaiming uniform; $\rho_{ij} = -5$, so

\begin{equation}
\sigma_{ij}\big|_{\rho=-5}
= \softplus(-5) = \ln(1 + e^{-5}) \approx 0.0067
\end{equation}

The network starts near-deterministic, giving the mean parameters time to reach
useful values before the posterior widens.

\subsection{Reparameterisation trick}

To allow gradient-based optimisation through stochastic weights:

\begin{equation}
w_{ij} = \mu_{ij} + \sigma_{ij}\,\varepsilon_{ij},
\qquad
\varepsilon_{ij} \sim \mathcal{N}(0, 1)
\label{eq:reparam}
\end{equation}

The randomness is moved into $\varepsilon_{ij}$, which carries no parameters.
Gradients flow through $\mu_{ij}$ and $\rho_{ij}$:

\begin{equation}
\frac{\partial w_{ij}}{\partial \mu_{ij}} = 1,
\qquad
\frac{\partial w_{ij}}{\partial \rho_{ij}}
= \varepsilon_{ij}\,\frac{d}{d\rho}\softplus(\rho_{ij})
= \varepsilon_{ij}\,\frac{e^{\rho_{ij}}}{1 + e^{\rho_{ij}}}
= \frac{\varepsilon_{ij}}{1 + e^{-\rho_{ij}}}
\end{equation}

The last form is the logistic sigmoid $s(\rho) = 1/(1+e^{-\rho})$ evaluated at $\rho_{ij}$.
Note that $\frac{d}{d\rho}\softplus(\rho) = s(\rho)$, \emph{not} $s(\rho)(1-s(\rho))$;
the latter is the self-derivative identity of the logistic sigmoid itself and does not apply here.

\subsection{KL divergence (\code{bayes.py:BayesianLinear.kl})}

The prior is $p(w) = \mathcal{N}(0,\, \sigma_{\mathrm{prior}}^2)$.
The closed-form KL between posterior and prior is:

\begin{equation}
\mathrm{KL}\!\bigl(q \,\|\, p\bigr)
= \sum_{i,j}
  \left[
    \ln\frac{\sigma_{\mathrm{prior}}}{\sigma_{ij}}
    + \frac{\sigma_{ij}^2 + \mu_{ij}^2}{2\,\sigma_{\mathrm{prior}}^2}
    - \frac{1}{2}
  \right]
\label{eq:kl}
\end{equation}

This equals zero when $\mu_{ij}=0$ and $\sigma_{ij}=\sigma_{\mathrm{prior}}$ (posterior
equals prior).  As the network learns, $\mu_{ij}$ shifts away from zero and
$\sigma_{ij}$ compresses, increasing the KL.  The term acts as a regulariser on
the magnitude and certainty of learned weights.

The total KL over the entire Bayesian MLP is the sum over all layers:

\begin{equation}
\mathrm{KL}_{\mathrm{total}} = \sum_{d=1}^{D} \mathrm{KL}_d
\end{equation}

\subsection{BayesianMLP output heads (\code{bayes.py:BayesianMLP})}

After $D$ Bayesian hidden layers with ReLU activations, two \emph{deterministic}
linear output heads produce:

\begin{align}
\hat{\mu} &= \mathbf{w}_\mu^\top\,\mathbf{x}^{(D)} + b_\mu \;\in \R
  \quad\text{(predicted affinity mean)}\\
s &= \mathrm{clamp}\!\bigl(
      \mathbf{w}_s^\top\,\mathbf{x}^{(D)} + b_s,\;
      {-10},\; 10
    \bigr) \;\in \R
  \quad\text{(log predictive variance, } s = \log\sigma_{\mathrm{al}}^2\text{)}
\end{align}

The clamp prevents numerical overflow in $e^{s}$ and $e^{-s}$.

\subsection{Heteroscedastic Gaussian NLL}

The model treats the affinity target as a Gaussian:
$y \sim \mathcal{N}(\hat{\mu},\, e^s)$.  The per-example negative
log-likelihood (dropping the additive constant $\tfrac{1}{2}\ln 2\pi$) is:

\begin{equation}
\ell_i = \frac{1}{2}\Bigl[e^{-s_i}\,(y_i - \hat{\mu}_i)^2 + s_i\Bigr]
\label{eq:hetero_nll}
\end{equation}

The term $e^{-s_i}$ reweights the squared error by the \emph{inverse predicted
variance}: when the model signals high aleatoric uncertainty (large $s_i$), the
squared error is discounted and the model only pays the entropy penalty $s_i$.
This makes the loss \emph{self-calibrating} — the model is penalised for being
overconfident and rewarded for expressing appropriate uncertainty.

The batch NLL is the mean:

\begin{equation}
\mathrm{NLL}_{\mathrm{batch}}
= \frac{1}{B} \sum_{i=1}^{B} \ell_i
\end{equation}

\subsection{ELBO training objective (\code{bayes.py:elbo\_loss})}

Variational inference maximises the Evidence Lower Bound~\cite{blundell2015}:

\begin{equation}
\mathcal{L}_{\mathrm{ELBO}}
= \mathbb{E}_{w \sim q}\!\bigl[\ln p(\mathbf{Y} \mid \mathbf{X}, w)\bigr]
  - \mathrm{KL}\!\bigl(q(w) \,\|\, p(w)\bigr)
\end{equation}

Minimising the negative ELBO per gradient step gives:

\begin{equation}
\boxed{
\mathcal{L}
= \underbrace{\mathrm{NLL}_{\mathrm{batch}}}_{\text{reconstruction}}
+ \underbrace{\beta \cdot \frac{\mathrm{KL}_{\mathrm{total}}}{N_{\mathrm{train}}}}_{\text{regularisation}}
}
\label{eq:elbo}
\end{equation}

where:
\begin{itemize}
  \item The NLL is estimated with \emph{one} weight sample per step (unbiased MC
        estimator of $\mathbb{E}_q[\mathrm{NLL}]$).
  \item $N_{\mathrm{train}}$ normalises the KL to per-example scale, preventing the
        regularisation term from dominating as the dataset grows.
  \item $\beta$ is the annealed KL weight (Section~\ref{sec:warmup}).
\end{itemize}

\subsection{KL warmup schedule (\code{train.py:train\_fusion\_from\_complexes})}
\label{sec:warmup}

\begin{equation}
\beta(\text{epoch}) = \beta_{\max} \cdot \min\!\left(1,\;
  \frac{\text{epoch}+1}{\tau}\right)
\label{eq:warmup}
\end{equation}

where $\beta_{\max}$ = \texttt{kl\_weight} (default 1.0) and $\tau$ =
\texttt{kl\_warmup} (default 10 epochs).

\medskip
\textbf{Why warmup?}  At epoch 0, a full KL penalty forces
$\sigma_{ij} \to \sigma_{\mathrm{prior}}$ (posterior collapses to the prior mean).
The means $\mu_{ij}$ receive no useful gradient, and the network never escapes the
prior.  By starting at $\beta = 0$ and ramping linearly, the network first
builds informative mean parameters before the regulariser becomes active.

% ============================================================
\section{Uncertainty Decomposition at Inference}
\label{sec:uncertainty}
% ============================================================

At test time, $T$ independent forward passes are drawn by sampling different
weight realisations:

\begin{equation}
\bigl\{(\hat{\mu}_t,\, s_t)\bigr\}_{t=1}^{T}, \quad
\hat{\mu}_t = \hat{\mu}(w_t),\quad w_t \sim q(w)
\end{equation}

\paragraph{Predictive mean.}
\begin{equation}
\bar{y} = \frac{1}{T}\sum_{t=1}^{T} \hat{\mu}_t
\end{equation}

\paragraph{Epistemic uncertainty.}  Model uncertainty, reducible with more data:
\begin{equation}
\sigma_{\mathrm{ep}}^2
= \mathrm{Var}_t(\hat{\mu}_t)
= \frac{1}{T}\sum_{t=1}^{T}(\hat{\mu}_t - \bar{y})^2
\label{eq:epistemic}
\end{equation}

\paragraph{Aleatoric uncertainty.}  Irreducible noise intrinsic to the data:
\begin{equation}
\sigma_{\mathrm{al}}^2
= \frac{1}{T}\sum_{t=1}^{T} e^{s_t}
= \mathbb{E}_t\!\bigl[e^{s_t}\bigr]
\label{eq:aleatoric}
\end{equation}

\paragraph{Total predictive variance.}  By the law of total variance:
\begin{equation}
\sigma_{\mathrm{total}}^2
= \underbrace{\mathrm{Var}_t(\hat{\mu}_t)}_{\sigma_{\mathrm{ep}}^2}
+ \underbrace{\mathbb{E}_t\!\bigl[\sigma_{\mathrm{al},t}^2\bigr]}_{\sigma_{\mathrm{al}}^2}
= \mathrm{Var}[Y] = \mathrm{Var}[\mathbb{E}[Y|W]] + \mathbb{E}[\mathrm{Var}[Y|W]]
\label{eq:totalvar}
\end{equation}

Both standard deviations ($\sigma_{\mathrm{ep}}$, $\sigma_{\mathrm{al}}$) are
returned in the same units as the affinity target (e.g.\ pK\textsubscript{d} units).

% ============================================================
\section{Training Losses for Standalone Models}
\label{sec:losses}
% ============================================================

The standalone GCN and GINE models (no Bayesian head) are trained with:

\begin{center}
\renewcommand{\arraystretch}{1.5}
\begin{tabular}{lll}
\toprule
\textbf{Task} & \textbf{Loss} & \textbf{Implemented in} \\
\midrule
Classification &
  $\mathcal{L}_{\mathrm{BCE}} = -\dfrac{1}{B}\displaystyle\sum_{i}
    \bigl[y_i\ln\sigma(\hat{y}_i) + (1-y_i)\ln(1-\sigma(\hat{y}_i))\bigr]$
  & \code{train.py:step} \\[10pt]
Regression &
  $\mathcal{L}_{\mathrm{MSE}} = \dfrac{1}{B}\displaystyle\sum_{i}(\hat{y}_i - y_i)^2$
  & \code{train.py:step} \\
\bottomrule
\end{tabular}
\end{center}

\medskip
Evaluation metrics:

\begin{align}
\mathrm{MAE} &= \frac{1}{N}\sum_{i=1}^N |\hat{y}_i - y_i| \\[4pt]
R^2 &= 1 - \frac{\sum_i (\hat{y}_i - y_i)^2}{\sum_i (y_i - \bar{y})^2} \\[4pt]
\mathrm{AUC\text{-}ROC} &= \Pr\!\bigl(\hat{p}_{\mathrm{pos}} > \hat{p}_{\mathrm{neg}}\bigr)
\end{align}

% ============================================================
\section{Architecture Summary}
\label{sec:summary}
% ============================================================

\begin{center}
\renewcommand{\arraystretch}{1.4}
\begin{tabular}{lllll}
\toprule
\textbf{Component} & \textbf{Input} & \textbf{Output} & \textbf{Default dim} & \textbf{Source} \\
\midrule
Atom embedding (Tower 2) & $z_i \in \mathbb{Z}$ & $\bh_i^{(0)} \in \R^C$ & $C=128$ & \code{models.py} \\
GINE layers (Tower 1)    & $\bx_i^{(0)} \in \R^7$ & $\bh_{\mathrm{lig}} \in \R^C$ & $L=3$ & \code{models.py} \\
GaussianRBF              & $d_{ij} \in \R$ & $\bphi(d_{ij}) \in \R^K$ & $K=50$ & \code{models.py} \\
InteractionBlock (×$T$)  & $\bh^{(\ell)}, \mathcal{E}, \bphi$ & $\bh^{(\ell+1)} \in \R^C$ & $T=3$ & \code{models.py} \\
Geometry readout         & $\bh_i^{(T)}$ & $\bh_{\mathrm{geo}} \in \R^C$ & $C=128$ & \code{models.py} \\
Fusion concat            & $\bh_{\mathrm{lig}}, \bh_{\mathrm{geo}}$ & $\mathbf{h} \in \R^{2C}$ & $2C=256$ & \code{models.py} \\
BayesianLinear (×$D$)    & $\mathbf{h}$ & $\mathbf{x}^{(D)} \in \R^{H_D}$ & $[256,128]$ & \code{bayes.py} \\
Output heads             & $\mathbf{x}^{(D)}$ & $\hat{\mu}, s \in \R$ & --- & \code{bayes.py} \\
MC uncertainty (×$T$)   & $T$ samples & $\bar{y}, \sigma_{\mathrm{ep}}, \sigma_{\mathrm{al}}$ & $T=30$ & \code{models.py} \\
\bottomrule
\end{tabular}
\end{center}

% ============================================================
\begin{thebibliography}{9}
\bibitem{kipf2017}
T.~N. Kipf and M.~Welling,
\textit{Semi-supervised classification with graph convolutional networks},
ICLR, 2017.

\bibitem{xu2019}
K.~Xu, W.~Hu, J.~Leskovec, and S.~Jegelka,
\textit{How powerful are graph neural networks?},
ICLR, 2019.

\bibitem{hu2020}
W.~Hu, B.~Liu, J.~Gomes, M.~Zitnik, P.~Liang, V.~Pande, and J.~Leskovec,
\textit{Strategies for pre-training graph neural networks},
ICLR, 2020.

\bibitem{schutt2017}
K.~T. Sch\"{u}tt, F.~Arbabzadah, S.~Chmiela, K.~R. M\"{u}ller, and
A.~Tkatchenko,
\textit{Quantum-chemical insights from deep tensor neural networks},
Nature Communications, 2017.

\bibitem{blundell2015}
C.~Blundell, J.~Cornebise, K.~Kavukcuoglu, and D.~Wierstra,
\textit{Weight uncertainty in neural networks},
ICML, 2015.
\end{thebibliography}

\end{document}
